\chapter{Quarks, QCD und Confinement}

\section{Ziel dieses Kapitels}

Dieses Kapitel behandelt die starke Wechselwirkung auf Quark- und Gluon-Ebene.
Im Mittelpunkt stehen Farbladung, Quantenchromodynamik, Gluonen,
Confinement, asymptotische Freiheit und Jets.

\begin{notebox}
Dieses Kapitel behandelt die Stichwoerter:
\begin{itemize}
    \item Quarks,
    \item Farbladung,
    \item Farbneutralitaet,
    \item Gluonen,
    \item Quantenchromodynamik,
    \item starke Kopplungskonstante,
    \item laufende Kopplung,
    \item asymptotische Freiheit,
    \item Confinement,
    \item Quark-Antiquark-Potential,
    \item String-Modell der Hadronen,
    \item Hadronisierung,
    \item Jets,
    \item experimenteller Hinweis auf Gluonen,
    \item Gluon-Selbstwechselwirkung.
\end{itemize}
\end{notebox}

\section{Warum braucht man QCD?}

Das Quarkmodell erklaert den Aufbau von Hadronen.
Protonen, Neutronen, Pionen, Kaonen und viele weitere Hadronen bestehen aus Quarks
und Antiquarks.

\begin{conceptbox}
Das Quarkmodell beantwortet die Frage:
\begin{center}
Welche Quarks stecken in einem Hadron?
\end{center}

Die Quantenchromodynamik beantwortet die Frage:
\begin{center}
Welche Wechselwirkung bindet Quarks zu Hadronen?
\end{center}
\end{conceptbox}

\begin{definitionbox}
Die \emph{Quantenchromodynamik}, kurz QCD, ist die Quantenfeldtheorie der starken
Wechselwirkung.

Sie beschreibt die Wechselwirkung von Quarks und Gluonen ueber Farbladung.
\end{definitionbox}

\begin{conceptbox}
Die QCD ist zur starken Wechselwirkung das, was die Quantenelektrodynamik zur
elektromagnetischen Wechselwirkung ist.

In der Elektrodynamik koppeln Photonen an elektrische Ladung.

In der QCD koppeln Gluonen an Farbladung.
\end{conceptbox}

\section{Farbladung}

Quarks tragen nicht nur elektrische Ladung, sondern auch Farbladung.
Der Begriff Farbe ist dabei nur ein Name.
Er hat nichts mit optischer Farbe zu tun.

\begin{definitionbox}
\emph{Farbladung} ist die Ladung der starken Wechselwirkung.
Quarks tragen Farbladung und wechselwirken daher stark.
\end{definitionbox}

\begin{conceptbox}
Es gibt drei Farbladungen, die oft bezeichnet werden als:
\[
    r,\qquad g,\qquad b.
\]
Diese Namen stehen fuer:
\[
    \text{rot},\qquad \text{gruen},\qquad \text{blau}.
\]

Antiquarks tragen entsprechende Antifarben:
\[
    \overline{r},
    \qquad
    \overline{g},
    \qquad
    \overline{b}.
\]
\end{conceptbox}

\begin{notebox}
Die Farbnamen sind nur eine Analogie.
Ein Quark ist nicht wirklich rot, gruen oder blau sichtbar.

Farbladung ist eine abstrakte Quantenzahl der starken Wechselwirkung.
\end{notebox}

\begin{examtipbox}
Elektrische Ladung gehoert zur elektromagnetischen Wechselwirkung.

Farbladung gehoert zur starken Wechselwirkung.
\end{examtipbox}

\section{Warum braucht man Farbe?}

Farbladung wurde eingefuehrt, um Hadronen und das Pauli-Prinzip konsistent zu
beschreiben.

\begin{conceptbox}
Baryonen bestehen aus drei Quarks.
Ein bekanntes Beispiel ist:
\[
    \Delta^{++}=uuu.
\]

Das $\Delta^{++}$ besitzt drei gleiche Up-Quarks.
Damit die Gesamtwellenfunktion mit dem Pauli-Prinzip vertraeglich ist, braucht man
eine zusaetzliche Quantenzahl.

Diese zusaetzliche Quantenzahl ist die Farbe.
\end{conceptbox}

\begin{conceptbox}
Die drei Up-Quarks im $\Delta^{++}$ koennen verschiedene Farben tragen:
\[
    u_r,\qquad u_g,\qquad u_b.
\]

Dadurch sind sie nicht in allen Quantenzahlen identisch.
Das erlaubt eine antisymmetrische Farbwellenfunktion.
\end{conceptbox}

\begin{examtipbox}
Ein wichtiger Grund fuer die Farbladung ist:
Sie macht Baryonen mit drei scheinbar gleichen Quarks mit dem Pauli-Prinzip
vertraeglich.
\end{examtipbox}

\section{Farbneutralitaet}

Einzelne Quarks werden nicht frei beobachtet.
Beobachtbare Hadronen sind farbneutral.

\begin{definitionbox}
Ein Zustand ist \emph{farbneutral}, wenn seine Gesamtfarbladung verschwindet.
Solche Zustaende nennt man auch Farbsinguletts.
\end{definitionbox}

\begin{conceptbox}
Baryonen bestehen aus drei Quarks mit drei verschiedenen Farben:
\[
    r+g+b
    \quad\Rightarrow\quad
    \text{farbneutral}.
\]

Mesonen bestehen aus einem Quark und einem Antiquark mit passender Farbe und
Antifarbe:
\[
    r+\overline{r}
    \quad\Rightarrow\quad
    \text{farbneutral}.
\]
\end{conceptbox}

\begin{formulabox}
Baryon:
\[
    q_r q_g q_b.
\]

Meson:
\[
    q_c \overline{q}_{\overline{c}}.
\]
\end{formulabox}

\begin{examtipbox}
Nur farbneutrale Zustaende werden als freie Teilchen beobachtet.

Quarks und Gluonen tragen Farbladung und treten nicht isoliert als freie Teilchen auf.
\end{examtipbox}

\section{Gluonen}

Die Austauschteilchen der starken Wechselwirkung sind Gluonen.

\begin{definitionbox}
\emph{Gluonen} sind die Austauschteilchen der starken Wechselwirkung.
Sie koppeln an Farbladung.
\end{definitionbox}

\begin{conceptbox}
Ein Quark kann ein Gluon emittieren oder absorbieren.
Dabei kann sich seine Farbe aendern, aber sein Flavor bleibt bei der starken
Wechselwirkung gleich.

Ein Up-Quark bleibt also ein Up-Quark.
Ein Strange-Quark bleibt ein Strange-Quark.
\end{conceptbox}

\begin{formulabox}
Typischer starker Vertex:
\[
    q \rightarrow q + g.
\]
\end{formulabox}

\begin{examplebox}
Ein rotes Quark kann durch Gluonemission in ein gruenes Quark uebergehen:
\[
    q_r \rightarrow q_g + g.
\]

Das Gluon traegt dann die passende Farbladung, sodass die Gesamtfarbladung erhalten
bleibt.
\end{examplebox}

\begin{examtipbox}
Gluonen aendern Farbe, aber nicht den Quark-Flavor.

Ein Flavorwechsel wie
\[
    s \rightarrow u
\]
ist nicht stark, sondern schwach.
\end{examtipbox}

\section{Acht Gluonen}

In der QCD gibt es nicht nur ein Gluon, sondern acht Gluonen.

\begin{conceptbox}
Naiv koennte man aus drei Farben und drei Antifarben neun Kombinationen bilden:
\[
    r\overline{r},
    \quad
    r\overline{g},
    \quad
    r\overline{b},
    \quad
    g\overline{r},
    \quad
    \ldots
\]

Die QCD besitzt aber die Symmetriegruppe
\[
    SU(3)_{\mathrm{Farbe}}.
\]
Daraus ergeben sich acht unabhaengige Gluon-Zustaende.
\end{conceptbox}

\begin{notebox}
Die Aussage, dass es acht Gluonen gibt, ist ein wichtiges Standardmodell-Ergebnis.

Fuer viele Aufgaben reicht:
\begin{center}
Gluonen tragen Farbladung und es gibt acht verschiedene Gluonfelder.
\end{center}
\end{notebox}

\begin{examtipbox}
Photon:
ein Austauschteilchen der elektromagnetischen Wechselwirkung.

Gluonen:
acht Austauschteilchen der starken Wechselwirkung.
\end{examtipbox}

\section{Photonen und Gluonen im Vergleich}

\begin{center}
\begin{tabular}{|p{4.0cm}|p{5.0cm}|p{5.0cm}|}
\hline
Eigenschaft & Photon & Gluon \\
\hline
Wechselwirkung & elektromagnetisch & stark \\
\hline
koppelt an & elektrische Ladung & Farbladung \\
\hline
eigene Ladung & elektrisch neutral & traegt Farbladung \\
\hline
Selbstwechselwirkung & keine direkte Photon-Photon-Kopplung in einfacher QED & direkte Gluon-Gluon-Kopplung \\
\hline
Reichweite & unendlich fuer freie elektromagnetische Wechselwirkung & effektiv auf Hadronen beschraenkt wegen Confinement \\
\hline
\end{tabular}
\end{center}

\begin{conceptbox}
Der wichtigste Unterschied ist:
Photonen tragen keine elektrische Ladung.
Gluonen tragen selbst Farbladung.

Deshalb wechselwirken Gluonen miteinander.
Diese Gluon-Selbstwechselwirkung ist zentral fuer die QCD.
\end{conceptbox}

\section{Gluon-Selbstwechselwirkung}

\begin{definitionbox}
\emph{Gluon-Selbstwechselwirkung} bedeutet, dass Gluonen direkt mit anderen Gluonen
wechselwirken koennen.
\end{definitionbox}

\begin{conceptbox}
Da Gluonen Farbladung tragen, koennen sie selbst Quelle der starken Wechselwirkung
sein.

Das ist anders als beim Photon:
Das Photon ist elektrisch neutral und koppelt nicht direkt an andere Photonen.
\end{conceptbox}

\begin{notebox}
In Feynman-Diagrammen hoeherer Ordnung der QCD treten daher Vertices mit mehreren
Gluonen auf.

Diese Selbstwechselwirkung ist wichtig fuer:
\begin{itemize}
    \item Confinement,
    \item asymptotische Freiheit,
    \item Jetbildung,
    \item nichtlineare Struktur der QCD.
\end{itemize}
\end{notebox}

\begin{examtipbox}
Gluon-Selbstwechselwirkung ist einer der zentralen Unterschiede zwischen QCD und QED.
\end{examtipbox}

\section{Starke Kopplungskonstante}

Die Staerke der starken Wechselwirkung wird durch eine Kopplung beschrieben.

\begin{definitionbox}
Die \emph{starke Kopplungskonstante} wird meist mit
\[
    \alpha_s
\]
bezeichnet.
Sie beschreibt die effektive Staerke der starken Wechselwirkung.
\end{definitionbox}

\begin{conceptbox}
Anders als eine einfache konstante Zahl haengt
\[
    \alpha_s
\]
von der Energieskala beziehungsweise vom Impulsuebertrag ab.

Man sagt:
Die Kopplung laeuft.
\end{conceptbox}

\begin{definitionbox}
Eine \emph{laufende Kopplung} ist eine Kopplungskonstante, deren effektiver Wert von
der betrachteten Energieskala abhaengt.
\end{definitionbox}

\begin{examtipbox}
Bei hoher Energie beziehungsweise kleinem Abstand wird die starke Kopplung kleiner.

Bei niedriger Energie beziehungsweise grossem Abstand wird sie groesser.
\end{examtipbox}

\section{Asymptotische Freiheit}

Eine der wichtigsten Eigenschaften der QCD ist asymptotische Freiheit.

\begin{definitionbox}
\emph{Asymptotische Freiheit} bedeutet:
Quarks wechselwirken bei sehr kleinen Abstaenden beziehungsweise sehr hohen
Impulsuebertraegen nur schwach miteinander.
\end{definitionbox}

\begin{conceptbox}
Bei grossem Impulsuebertrag gilt:
\[
    Q^2 \text{ gross}
    \quad\Rightarrow\quad
    \alpha_s(Q^2) \text{ klein}.
\]

Dann koennen Quarks innerhalb eines Hadrons naeherungsweise wie fast freie Teilchen
erscheinen.
\end{conceptbox}

\begin{notebox}
Das ist wichtig fuer tiefinelastische Streuung.
Bei hohen Impulsuebertraegen sieht man punktfoermige Bestandteile im Proton.
Diese Bestandteile werden als Quarks beziehungsweise Partonen interpretiert.
\end{notebox}

\begin{examtipbox}
Asymptotische Freiheit:
\[
    \text{kleiner Abstand}
    \quad\Rightarrow\quad
    \text{schwache effektive starke Kopplung}.
\]

Das klingt zuerst paradox, ist aber eine zentrale Eigenschaft der QCD.
\end{examtipbox}

\section{Confinement}

Die zweite zentrale Eigenschaft der QCD ist Confinement.

\begin{definitionbox}
\emph{Confinement} bedeutet:
Quarks und Gluonen werden nicht als freie isolierte Teilchen beobachtet.
Sie sind in farbneutralen Hadronen eingeschlossen.
\end{definitionbox}

\begin{conceptbox}
Wenn man versucht, zwei Quarks voneinander zu trennen, wird die starke Feldenergie
zwischen ihnen groesser.

Statt ein einzelnes freies Quark zu erhalten, wird bei genuegend grosser Energie ein
neues Quark-Antiquark-Paar erzeugt.

Dadurch entstehen neue Hadronen.
\end{conceptbox}

\begin{examtipbox}
Confinement erklaert, warum man keine isolierten Quarks im Detektor sieht.

Man sieht stattdessen Hadronen oder Jets aus Hadronen.
\end{examtipbox}

\section{Quark-Antiquark-Potential}

Das Potential zwischen einem Quark und einem Antiquark hat bei grossen Abstaenden
einen naeherungsweise linearen Anteil.

\begin{formulabox}
Eine haeufige qualitative Form ist:
\[
    V(r)
    =
    -\frac{a}{r}
    +
    \kappa r.
\]
Dabei beschreibt:
\begin{itemize}
    \item $-\frac{a}{r}$ einen coulombartigen Anteil bei kleinen Abstaenden,
    \item $\kappa r$ einen linearen Confinement-Anteil bei grossen Abstaenden.
\end{itemize}
\end{formulabox}

\begin{conceptbox}
Bei kleinen Abstaenden dominiert der coulombartige Anteil.
Das passt zur asymptotischen Freiheit.

Bei grossen Abstaenden waechst das Potential ungefaehr linear.
Das bedeutet:
Je weiter man Quarks trennt, desto mehr Energie steckt im Farbfeld.
\end{conceptbox}

\begin{examtipbox}
Quark-Antiquark-Potential:
\[
    \text{kleines } r:
    \quad
    \text{coulombartig}.
\]
\[
    \text{grosses } r:
    \quad
    \text{linear ansteigend}.
\]

Der lineare Anteil ist der anschauliche Ursprung des Confinements.
\end{examtipbox}

\section{String-Modell der Hadronen}

Das lineare Potential fuehrt zur Vorstellung eines Farbfeld-Strings zwischen Quark
und Antiquark.

\begin{definitionbox}
Im \emph{String-Modell} der Hadronen wird das Farbfeld zwischen Quark und Antiquark
als schlauchartiger String beschrieben.
\end{definitionbox}

\begin{conceptbox}
Wenn Quark und Antiquark auseinandergezogen werden, wird der String laenger.
Die im String gespeicherte Energie waechst.

Irgendwann ist genug Energie vorhanden, um ein neues Quark-Antiquark-Paar zu erzeugen:
\[
    q\overline{q}.
\]

Dann reisst der String nicht in freie Quarks, sondern es entstehen zwei farbneutrale
Hadronen.
\end{conceptbox}

\begin{formulabox}
Schematisch:
\[
    q\overline{q}
    \rightarrow
    q\overline{q}
    +
    q\overline{q}.
\]
\end{formulabox}

\begin{examtipbox}
Beim Auseinanderziehen eines Quark-Antiquark-Paares entstehen neue Hadronen.
Man bekommt keine freien Quarks.
\end{examtipbox}

\section{Hadronisierung}

Quarks und Gluonen koennen bei hohen Energien erzeugt werden.
Wegen Confinement erscheinen sie aber nicht frei im Detektor.

\begin{definitionbox}
\emph{Hadronisierung} ist der Prozess, bei dem farbige Quarks und Gluonen in
farbneutrale Hadronen uebergehen.
\end{definitionbox}

\begin{conceptbox}
Ein hochenergetisches Quark erzeugt durch Gluonabstrahlung und Paarerzeugung viele
weitere farbige Objekte.

Diese ordnen sich anschliessend zu farbneutralen Hadronen.
Das beobachtbare Ergebnis ist ein Buendel vieler Hadronen.
\end{conceptbox}

\begin{examtipbox}
Quarks und Gluonen werden nicht direkt gemessen.
Man rekonstruiert sie indirekt aus den Hadronen, die bei der Hadronisierung entstehen.
\end{examtipbox}

\section{Jets}

Ein Jet ist ein kollimiertes Buendel von Hadronen, das aus einem hochenergetischen
Quark oder Gluon entsteht.

\begin{definitionbox}
Ein \emph{Jet} ist ein gebuendelter Teilchenschauer aus Hadronen, der die Richtung
eines urspruenglichen hochenergetischen Quarks oder Gluons widerspiegelt.
\end{definitionbox}

\begin{conceptbox}
Wenn in einer hochenergetischen Kollision ein Quark erzeugt wird, kann es wegen
Confinement nicht frei auslaufen.

Stattdessen hadronisiert es.
Die entstehenden Hadronen fliegen ungefaehr in Richtung des urspruenglichen Quarks.
Das beobachtet man als Jet.
\end{conceptbox}

\begin{notebox}
Jets sind eines der wichtigsten experimentellen Signale fuer Quarks und Gluonen.

Man sieht nicht das einzelne Quark, sondern ein Buendel von Hadronen.
\end{notebox}

\begin{examtipbox}
Jet:
\[
    \text{experimentelles Signal eines hochenergetischen Quarks oder Gluons}.
\]

Confinement:
\[
    \text{Grund, warum man statt freier Quarks Jets sieht}.
\]
\end{examtipbox}

\section{Zwei-Jet-Ereignisse}

Ein einfaches Beispiel ist Elektron-Positron-Annihilation in ein Quark-Antiquark-Paar.

\begin{formulabox}
\[
    e^- + e^+
    \rightarrow
    \gamma^*
    \rightarrow
    q+\overline{q}.
\]
\end{formulabox}

\begin{conceptbox}
Das Quark und das Antiquark fliegen im Schwerpunktssystem ungefaehr in entgegengesetzte
Richtungen.

Beide hadronisieren.

Das Ergebnis sind zwei entgegengesetzte Jets:
\[
    e^-+e^+
    \rightarrow
    \text{Jet}+\text{Jet}.
\]
\end{conceptbox}

\begin{examtipbox}
Zwei-Jet-Ereignisse sind ein Hinweis auf die Erzeugung eines Quark-Antiquark-Paares.
\end{examtipbox}

\section{Drei-Jet-Ereignisse und Gluonnachweis}

Ein wichtiger experimenteller Hinweis auf Gluonen sind Drei-Jet-Ereignisse.

\begin{formulabox}
\[
    e^- + e^+
    \rightarrow
    q+\overline{q}+g.
\]
\end{formulabox}

\begin{conceptbox}
Wenn ein Quark oder Antiquark ein hartes Gluon abstrahlt, kann auch das Gluon
hadronisieren.

Dann entstehen drei Jets:
\[
    q \rightarrow \text{Jet},
\]
\[
    \overline{q} \rightarrow \text{Jet},
\]
\[
    g \rightarrow \text{Jet}.
\]
\end{conceptbox}

\begin{notebox}
Drei-Jet-Ereignisse in Elektron-Positron-Kollisionen waren ein wichtiger
experimenteller Hinweis auf die Existenz von Gluonen.
\end{notebox}

\begin{examtipbox}
Zwei Jets:
\[
    q\overline{q}.
\]

Drei Jets:
\[
    q\overline{q}g.
\]

Ein dritter Jet kann ein Hinweis auf ein abgestrahltes Gluon sein.
\end{examtipbox}

\section{Quarks als Partonen}

Bei hohen Impulsuebertraegen kann man Hadronen als aus punktfoermigen Bestandteilen
aufgebaut betrachten.
Diese Bestandteile nennt man Partonen.

\begin{definitionbox}
\emph{Partonen} sind die punktfoermigen Bestandteile eines Hadrons, die in
hochenergetischen Streuprozessen aufgeloest werden.
\end{definitionbox}

\begin{conceptbox}
Im modernen Bild bestehen Hadronen aus:
\begin{itemize}
    \item Valenzquarks,
    \item Seequarks,
    \item Gluonen.
\end{itemize}
\end{conceptbox}

\begin{definitionbox}
\emph{Valenzquarks} bestimmen die grundlegenden Quantenzahlen eines Hadrons.

\emph{Seequarks} sind Quark-Antiquark-Paare, die durch Quantenfluktuationen entstehen.

\emph{Gluonen} tragen einen Teil des Impulses und vermitteln die starke Wechselwirkung.
\end{definitionbox}

\begin{examplebox}
Das Proton hat als Valenzquarkinhalt:
\[
    p=uud.
\]

In einem hochenergetischen Streuexperiment enthaelt das Proton aber zusaetzlich
Seequarks und Gluonen.
\end{examplebox}

\section{Tiefinelastische Streuung und Substruktur}

Die tiefinelastische Elektron-Proton-Streuung zeigte, dass Protonen innere
punktfoermige Bestandteile besitzen.

\begin{conceptbox}
Bei elastischer Streuung bleibt das Proton als Ganzes erhalten.

Bei tiefinelastischer Streuung wird das Proton aufgebrochen.
Das Elektron streut an einem einzelnen Parton im Proton.
\end{conceptbox}

\begin{notebox}
Ein grosser Impulsuebertrag bedeutet kleine raeumliche Aufloesung.
Dadurch koennen die inneren Bestandteile des Protons sichtbar werden.
\end{notebox}

\begin{examtipbox}
Experimenteller Nachweis der Nukleonensubstruktur:
hochenergetische inelastische Elektron-Proton-Streuung.

Interpretation:
Das Proton enthaelt punktfoermige Partonen, insbesondere Quarks.
\end{examtipbox}

\section{Farbe und Hadronen}

\begin{conceptbox}
Alle beobachtbaren Hadronen sind farbneutral.

Baryonen:
\[
    q_r q_g q_b.
\]

Mesonen:
\[
    q_c\overline{q}_{\overline{c}}.
\]

Gluonreiche Zustaende oder exotische Hadronen muessen ebenfalls insgesamt
farbneutral sein.
\end{conceptbox}

\begin{notebox}
Auch exotische Hadronen wie Tetraquarks oder Pentaquarks sind nur beobachtbar,
wenn sie insgesamt farbneutral sind.
\end{notebox}

\section{Exotische Hadronen}

Das einfache Quarkmodell enthaelt Baryonen und Mesonen.
Die QCD erlaubt aber allgemein alle farbneutralen Kombinationen.

\begin{definitionbox}
\emph{Exotische Hadronen} sind Hadronen, die nicht einfach als normales Meson
\[
    q\overline{q}
\]
oder normales Baryon
\[
    qqq
\]
beschrieben werden.
\end{definitionbox}

\begin{conceptbox}
Moegliche exotische Hadronen sind:
\begin{itemize}
    \item Tetraquarks,
    \item Pentaquarks,
    \item Glueballs,
    \item hybride Mesonen.
\end{itemize}
\end{conceptbox}

\begin{notebox}
Fuer diese Vorlesung ist meistens die Grundidee wichtiger als Details:
Die QCD erlaubt farbneutrale gebundene Zustaende, solange die Farbladung insgesamt
verschwindet.
\end{notebox}

\section{QCD bei kleinen und grossen Abstaenden}

\begin{center}
\begin{tabular}{|p{4.2cm}|p{5.0cm}|p{5.0cm}|}
\hline
Bereich & Verhalten der QCD & physikalische Folge \\
\hline
kleiner Abstand & kleine effektive Kopplung & asymptotische Freiheit \\
\hline
grosser Abstand & grosse effektive Kopplung & Confinement \\
\hline
hohe Energie & perturbative Rechnungen oft moeglich & Quarks und Gluonen als Partonen \\
\hline
niedrige Energie & nicht-perturbativ & Hadronen, Bindung, Massenspektrum \\
\hline
\end{tabular}
\end{center}

\begin{examtipbox}
QCD ist bei hohen Energien oft einfacher zu rechnen als bei niedrigen Energien.

Grund:
Bei hohen Energien ist
\[
    \alpha_s
\]
kleiner.
\end{examtipbox}

\section{Starke und elektromagnetische Kopplung im Vergleich}

\begin{conceptbox}
Die elektromagnetische Kopplung ist durch die Feinstrukturkonstante
\[
    \alpha
\]
charakterisiert.

Die starke Kopplung wird durch
\[
    \alpha_s
\]
beschrieben.
\end{conceptbox}

\begin{notebox}
Die elektromagnetische Kopplung ist bei typischen Teilchenphysik-Energien relativ
klein.
Deshalb ist Stoerungsrechnung in der QED sehr erfolgreich.

Die starke Kopplung kann bei niedrigen Energien gross werden.
Dann sind einfache Stoerungsrechnungen nicht mehr ausreichend.
\end{notebox}

\begin{examtipbox}
QED:
kleine Kopplung, Photonen ohne Selbstwechselwirkung.

QCD:
laufende Kopplung, Gluon-Selbstwechselwirkung, Confinement.
\end{examtipbox}

\section{Mini-Zusammenfassung}

\begin{notebox}
Die wichtigsten Punkte dieses Kapitels sind:
\begin{itemize}
    \item QCD ist die Theorie der starken Wechselwirkung.
    \item Quarks tragen Farbladung.
    \item Gluonen vermitteln die starke Wechselwirkung.
    \item Gluonen tragen selbst Farbladung.
    \item Deshalb gibt es Gluon-Selbstwechselwirkung.
    \item Beobachtbare Hadronen sind farbneutral.
    \item Baryonen sind farbneutral als Kombination von drei Farben.
    \item Mesonen sind farbneutral als Farbe plus Antifarbe.
    \item Die starke Kopplung
    \[
        \alpha_s
    \]
    haengt von der Energieskala ab.
    \item Asymptotische Freiheit bedeutet schwache Kopplung bei kleinen Abstaenden.
    \item Confinement bedeutet, dass Quarks und Gluonen nicht frei beobachtet werden.
    \item Das Quark-Antiquark-Potential hat qualitativ die Form:
    \[
        V(r)=-\frac{a}{r}+\kappa r.
    \]
    \item Das String-Modell erklaert anschaulich Hadronisierung.
    \item Jets sind experimentelle Signale von hochenergetischen Quarks oder Gluonen.
    \item Drei-Jet-Ereignisse sind ein Hinweis auf Gluonen.
\end{itemize}
\end{notebox}

\section{Pruefungscheck}

\begin{examtipbox}
Nach diesem Kapitel solltest du folgende Fragen beantworten koennen:
\begin{enumerate}
    \item Was beschreibt die QCD?
    \item Was ist Farbladung?
    \item Warum ist Farbe keine optische Farbe?
    \item Warum braucht man Farbladung im Quarkmodell?
    \item Was bedeutet Farbneutralitaet?
    \item Wie werden Baryonen farbneutral?
    \item Wie werden Mesonen farbneutral?
    \item Was sind Gluonen?
    \item Warum gibt es Gluon-Selbstwechselwirkung?
    \item Was ist der Unterschied zwischen Photon und Gluon?
    \item Was bedeutet laufende Kopplung?
    \item Was ist asymptotische Freiheit?
    \item Was ist Confinement?
    \item Warum sieht man keine freien Quarks?
    \item Welche qualitative Form hat das Quark-Antiquark-Potential?
    \item Was beschreibt das String-Modell der Hadronen?
    \item Was ist Hadronisierung?
    \item Was ist ein Jet?
    \item Warum sind Jets Hinweise auf Quarks und Gluonen?
    \item Warum sind Drei-Jet-Ereignisse Hinweise auf Gluonen?
    \item Was sind Valenzquarks, Seequarks und Gluonen im Proton?
    \item Was zeigt tiefinelastische Streuung ueber die Nukleonensubstruktur?
\end{enumerate}
\end{examtipbox}